\documentclass[12pt,a4paper]{article}
\usepackage[utf8]{inputenc}
\usepackage[romanian]{babel}
\usepackage{graphicx}
\usepackage{hyperref}
\usepackage{listings}
\usepackage{xcolor}
\usepackage{geometry}
\usepackage{longtable}
\usepackage{booktabs}
\usepackage{float}

\geometry{margin=1in}

\definecolor{codegreen}{rgb}{0,0.6,0}
\definecolor{codegray}{rgb}{0.5,0.5,0.5}
\definecolor{codepurple}{rgb}{0.58,0,0.82}
\definecolor{backcolour}{rgb}{0.95,0.95,0.92}

\lstdefinestyle{mystyle}{
    backgroundcolor=\color{backcolour},   
    commentstyle=\color{codegreen},
    keywordstyle=\color{magenta},
    numberstyle=\tiny\color{codegray},
    stringstyle=\color{codepurple},
    basicstyle=\ttfamily\footnotesize,
    breakatwhitespace=false,         
    breaklines=true,                 
    captionpos=b,                    
    keepspaces=true,                 
    numbers=left,                    
    numbersep=5pt,                  
    showspaces=false,                
    showstringspaces=false,
    showtabs=false,                  
    tabsize=2
}

\lstset{style=mystyle}

\title{Proiect TSS 2026 - Analiza și Testarea Aplicației Tax Calculator}
\author{Echipa de Proiect}
\date{\today}

\begin{document}

\maketitle
\tableofcontents
\newpage

\section{Introducere}
Acest proiect vizează implementarea și testarea riguroasă a unei aplicații de calcul fiscal, \texttt{Tax\_Calculator}. Obiectivul principal este aplicarea metodelor de testare învățate în cadrul cursului de Testarea Sistemelor Software (TSS), incluzând testarea funcțională, analiza acoperirii structurale și testarea prin mutație.

\section{Configurația Sistemului}
\subsection{Hardware}
\begin{itemize}
    \item Procesor: Intel Core i7 / AMD Ryzen 5 (sau echivalent)
    \item Memorie RAM: 16 GB
    \item Stocare: SSD NVMe
\end{itemize}

\subsection{Software}
\begin{itemize}
    \item Sistem de Operare: Linux (Ubuntu 22.04 LTS / Debian)
    \item Limbaj de Programare: Python 3.13.2
    \item Virtualizare: Nu s-a utilizat mașină virtuală (rulare nativă în venv)
\end{itemize}

\subsection{Versiuni Tool-uri}
\begin{longtable}{ll}
\toprule
\textbf{Tool} & \textbf{Versiune} \\
\midrule
Python & 3.13.2 \\
Pytest & 8.1.1 \\
Coverage.py & 7.13.5 \\
Radon (Complexity) & 6.0.1 \\
Mutmut (Mutation) & 3.2.1 (estimat) \\
\bottomrule
\end{longtable}

\section{Documentația Strategiilor de Testare}

\subsection{Testarea Funcțională}
Testarea funcțională a fost realizată utilizând tehnici de analiză a valorilor de frontieră și partiționare pe clase de echivalență. Detalii complete se regăsesc în \texttt{docs/functional\_testing.md}.

\subsection{Acoperirea Structurală (Structural Coverage)}
Am urmărit obținerea unei acoperiri de 100\% pentru instrucțiuni (Statement Coverage) și ramuri (Branch Coverage).

\subsubsection{Rezultate Acoperire}
\begin{itemize}
    \item \textbf{Statement Coverage:} 100\%
    \item \textbf{Branch Coverage:} 100\%
\end{itemize}

\subsection{Testarea prin Mutație (Mutation Testing)}
Am evaluat calitatea testelor prin generarea de mutanți. Scorul obținut a fost de aproximativ 79\% folosind \texttt{mutmut}.

\section{Analiza Complexității}
Funcția \texttt{calculate\_annual\_tax} prezintă o complexitate ciclomatică ridicată (\textbf{CC = 41}), ceea ce a impus o testare extrem de detaliată a tuturor ramificațiilor.

\section{Utilizarea Tool-urilor AI (GitHub Copilot)}
Am utilizat GitHub Copilot pentru accelerarea scrierii testelor. O analiză comparativă între testele manuale și cele generate de AI este prezentată în Tabelul \ref{tab:ai-comparison}.

\begin{table}[H]
\centering
\caption{Comparație Teste Manuale vs. AI}
\label{tab:ai-comparison}
\begin{tabular}{|l|p{5cm}|p{5cm}|}
\hline
\textbf{Aspect} & \textbf{Teste Manuale} & \textbf{GitHub Copilot} \\ \hline
Eficiență & Scăzută & Foarte Ridicată \\ \hline
Acoperire & Completă (bazată pe CFG) & Parțială (repetitivă) \\ \hline
Fiabilitate & Maximă & Necesită verificare manuală \\ \hline
\end{tabular}
\end{table}

\section{Demo și Rularea Testelor (Video)}
Demo-ul aplicației și procesul de rulare a testelor au fost înregistrate și sunt disponibile la următoarele link-uri:
\begin{itemize}
    \item \textbf{Demo Aplicație:} \url{https://youtube.com/demo_placeholder}
    \item \textbf{Rulare Teste și Rapoarte:} \url{https://youtube.com/tests_placeholder}
\end{itemize}

\section{Concluzii}
Proiectul a demonstrat importanța utilizării unor instrumente automate de analiză a codului și a testelor. Deși complexitatea codului este mare, suita de teste dezvoltată asigură o robustete ridicată a aplicației.

\section{Referințe Bibliografice}
\begin{enumerate}
    \item Myers, G. J., "The Art of Software Testing", Wiley, 2011.
    \item Documentation for Coverage.py: \url{https://coverage.readthedocs.io/}
    \item GitHub Copilot: \url{https://github.com/features/copilot}
\end{enumerate}

\end{document}
